\subsubsection*{恒力功、功率与动能定理的使用范围}

恒力的大小和方向都不变，可从式~\eqref{ch08:eq:work-integral} 的积分号内提出：
\[
 W=\vec F\cdot\int_A^B\dd\vec r
  =\vec F\cdot(\vec r_B-\vec r_A)=F|\vec r_B-\vec r_A|\cos\theta.
\]
这里用的是始末位置之间的位移，不是弯曲路径的长度。变力一般不能这样提出积分号，必须沿实际路径计算。
点积还说明只有切向分量作功；法向力即使很大也可能不作功，但它仍可改变速度方向。

功率表示作功的快慢。时间间隔 $\Delta t$ 内的平均功率为 $\overline P=W/\Delta t$；令时间间隔趋于零，得到瞬时功率
\begin{equation}
 P=\frac{\dd W}{\dd t}
 =\vec F\cdot\frac{\dd\vec r}{\dd t}
 =\vec F\cdot\vec v=F_t v.
 \label{ch08:eq:power}
\end{equation}
$F_t$ 是力沿速度方向的有符号分量。功的单位是焦耳，$1\,\mathrm J=1\,\mathrm{N\,m}$；功率单位是瓦特，
$1\,\mathrm W=1\,\mathrm{J/s}$。只有牵引力平行于速度且 $v>0$ 时，才能由此写成 $F=P/v$。
因此匀功率与恒牵引力不是同一条件：匀功率时速度升高，牵引力减小。不能把这个理想关系直接外推到静止启动瞬间；实际机器的牵引力有上限。

也可以在直角坐标中逐项验证式~\eqref{ch08:eq:kinetic-rate}。把牛顿第二定律的三个分量分别乘以相应速度分量，再相加：
\begin{align*}
 F_xv_x+F_yv_y+F_zv_z
 &=m\left(v_x\frac{\dd v_x}{\dd t}+v_y\frac{\dd v_y}{\dd t}+v_z\frac{\dd v_z}{\dd t}\right)\\
 &=\frac{\dd}{\dd t}\left[\frac m2(v_x^2+v_y^2+v_z^2)\right]
 =\frac{\dd E_k}{\dd t}.
\end{align*}
再乘 $\dd t$ 并积分，左边利用 $v_x\dd t=\dd x$ 等关系化为合力功，右边成为末动能减初动能。
所以动能定理是牛顿第二定律沿实际路径的积分形式，适于联系速率和位置；若还要求时间或速度方向，通常要配合运动学或分量方程。

对质点系，外力功与内力功共同决定总动能变化。两个物体间作用力等大反向，但其作用点位移一般不同，
所以内力总功未必为零；弹簧伸缩、碰撞形变和摩擦发热都体现了这一点。只有把内力功转化为相应势能或内能，
才能写成更完整的能量关系。参考系改变时速度、动能和功都会改变，但在同一惯性系中二者始终满足动能定理。

虚功方法研究满足约束的假想无穷小位移。理想约束力对允许的虚位移的总功为零，平衡时主动外力的总虚功也为零。
它把合力、合力矩平衡改写到约束允许的坐标中，常能避开未知支持力和绳张力；后文将用动滑轮说明。
